\scanpage{296}
\mspage{246}

On a prouvé la

\uline{Proposition 1}8 Soit $S$ dans $M$. Il existe, à isom unique près, un foncteur et un seul
\[
 Ho^{\unsure{lc|le}}_{W\deleted{\unsure{,}}}(S)\times Ho^{lc}_W(S)\xrightarrow{\pi}Ho_{W,S}(\Phi)
\]
\deleted{\unsure{dont le composé avec}} \inserted{donnant lieu à un carré commutatif} à
\[
 \deleted{\underline{Fib}_W\,S\times\underline{Fib}_W\,S\longrightarrow Ho^{lc}_W(S)\times Ho^{lc}_W(S)}
\]
isom près
\diagram{\begin{tikzcd}[ampersand replacement=\&,column sep=2pt,row sep=tiny,font=\scriptsize]
 \& \& X_1 \& \\
 \& (X_1,X_2)\arrow[r,mapsto] \& {} \& X_1\times_S X_2\arrow[ul]\arrow[dl] \\
 \& \& X_2 \& \\
 \underline{Fib}_W\,S\times\underline{Fib}_W\,S\arrow[rrr]\arrow[d] \& \& \& \deleted{\unsure{\underline{Fib}_W}}M_{/S}(\Phi)\arrow[d]\\
 Ho^{lc}_W(S)\times Ho^{lc}_W(S)\arrow[rrr] \& \& \& Ho_{W,S}(\Phi)\\
 \bigl(=(M_{/S})(W_S)^{-1}\times(M_{/S})(W_S)^{-1} \& \& \& \bigl(=M_{/S}(\Phi)(W_S(\Phi))^{-1}\bigr)
\end{tikzcd}}
\deleted{\unsure{Consid.}} Pour $X_1,X_2$ dans $M_{/S}$, \deleted{\unsure{leur}} le diagramme
\[
 \pi(X_1,X_2)\overset{\mathrm{def}}{=}\pi\bigl(ho^{lc}_{W,S}(X_1),ho^{lc}_{W,S}(X_2)\bigr)
\]
se ``calcule'' ainsi à iso. can. près : on choisit
\[
 X_1\overset{\alpha_1}{\unsure{-|\to}}\overline X_1,\quad
 X_2\overset{\alpha_2}{-}\overline X_2\qquad\text{flèches dans }M_{/S}
\]
avec
\[
 \alpha_1,\alpha_2\in W_S,\qquad\overline X_1,\overline X_2\in\mathrm{Ob}\,\underline{Fib}_W\,S,
\]
ce qui définit
\diagram{\begin{tikzcd}[ampersand replacement=\&,column sep=small,row sep=small]
 \pi(\alpha_1,\alpha_2):\quad \pi(X_1,X_2)\arrow[r,"\sim"] \&
 \pi(\overline X_1,\overline X_2)=ho^{\unsure{\xi|\zeta}}_{\unsure{S,M,\Phi|S,W,\Phi}}\bigg( \&
 \overline X_1 \& \overline X_1\times_S\overline X_2\arrow[l]\arrow[d] \& \bigg)\\
 \& \& \&\overline X_2\& .
\end{tikzcd}}
% archive sheet number: 296

